\documentclass[11pt]{article}
\usepackage[margin=1in]{geometry}
\usepackage{amsmath,amssymb,booktabs,hyperref,graphicx}
\title{Replication Report: Three-dimensional Topological Orbital Hall Effect\\
Caused by Magnetic Hopfions}
\author{Replication package for G\"obel \& Lounis, arXiv:2506.11448v3 (2025)\\
Verdict: \textbf{PARTIAL}}
\date{Reproduced 2026-07-19}
\begin{document}
\maketitle

\begin{abstract}
G\"obel and Lounis predict a genuinely three-dimensional \emph{orbital} Hall effect
generated by a magnetic hopfion, with finite out-of-plane and in-plane orbital Hall
conductivities $\sigma^{L_z}$ at the Fermi level for exchange $m=7t$ and \emph{no}
spin--orbit coupling. We rebuild the non-collinear tight-binding + Kubo pipeline from
scratch on a $12\times12\times6$ cubic lattice, implement a hopfion (twisted bimeron
tube) texture, and evaluate $\sigma^{L_z}_{xy}$ via a real-space Kubo--Bastin formula
using the itinerant orbital operator $L_z=\tfrac12(X v_y - Y v_x)$. We confirm the
central qualitative claim---a finite orbital Hall conductivity without SOC, isolated
against a filling-matched uniform-FM reference. The strict \emph{quantitative}
topological separation, however, is a difference of two large numbers
($\sigma_{\text{hopfion}}=4330$, $\sigma_{\text{FM}}=4055$, residual $275\,[e/2\pi]$)
and is a known hard limit of the itinerant-$L_z$ operator under periodic boundary
conditions, shared with the sibling 2D paper (gobel2024). Hence the honest verdict is
\textbf{PARTIAL}.
\end{abstract}

\section{The paper's claim}
A magnetic hopfion is a bimeron/skyrmion tube that is twisted once and closed into a
torus; it carries an integer Hopf index (a 3D topological invariant). Unlike a 2D
skyrmion, whose 2D winding drives a \emph{topological Hall} effect, the hopfion is
argued to have no sizeable topological Hall signal but instead a unique 3D
\emph{orbital} Hall response. The local emergent field
\begin{equation}
B_{\text{em},\alpha}(\mathbf r)=\tfrac12\,\epsilon_{\alpha\beta\gamma}\,
\mathbf m(\mathbf r)\cdot\!\left(\frac{\partial \mathbf m}{\partial\beta}\times
\frac{\partial \mathbf m}{\partial\gamma}\right)
\end{equation}
seeds an orbital Berry curvature that produces both out-of-plane ($\sigma^{L_z}_{xy}$)
and in-plane ($\sigma^{L_x},\sigma^{L_y}$) orbital Hall conductivities, finite at the
Fermi level even in the absence of spin--orbit coupling. This orbital signature is
proposed as an electronic hallmark that distinguishes a hopfion from a look-alike
skyrmionium in real-space microscopy.

\section{Model built}
We reproduce the itinerant s--d tight-binding Hamiltonian
\begin{equation}
H = -t\sum_{\langle ij\rangle,\sigma} c^\dagger_{i\sigma}c_{j\sigma}
    + m\sum_i c^\dagger_i\,(\mathbf n_i\cdot\boldsymbol\sigma)\,c_i ,
\end{equation}
with $t=1$ (energy unit), exchange $m=7t$, on a cubic lattice of $L_x\times L_y\times
L_z = 12\times12\times6=864$ sites (spinful dimension $1728$). The texture
$\mathbf n_i=\mathbf m(\mathbf r_i)$ is a hopfion of size $\lambda=8a$: inside the
torus the in-plane bimeron angle twists helically by $2\pi$ along $z$ (the Hopf link
structure), and outside the texture relaxes to the uniform background $+\hat z$, giving
the periodicity required by the reciprocal-space picture (paper Eq.~1, unit cell
$2\lambda\times2\lambda\times\lambda$). No spin--orbit coupling is included anywhere;
the only source of orbital response is the real-space texture.

\section{Method built (orbital Hall via Kubo--Bastin)}
The itinerant orbital angular momentum along $z$ is
$L_z=\tfrac12(X_c v_y - Y_c v_x)$ (symmetrized, with center-shifted position
operators $X_c=X-c_x$, $Y_c=Y-c_y$), where velocities follow from the position
operators, $v_\alpha = i[H,R_\alpha]$. The orbital current is
$j^{L_z}_x=\tfrac12\{L_z,v_x\}$. The zero-temperature, DC, clean-limit Hall response is
the Kubo--Bastin (Kubo--Greenwood off-diagonal) sum
\begin{equation}
\sigma^{L_z}_{xy}
= i\sum_{n\in\text{occ}}\sum_{m\in\text{unocc}}
\frac{\langle n|j^{L_z}_x|m\rangle\langle m|v_y|n\rangle
     -\langle n|v_y|m\rangle\langle m|j^{L_z}_x|n\rangle}{(E_n-E_m)^2},
\end{equation}
evaluated in the real-space ($\Gamma$-point) supercell. States degenerate across the
Fermi level (denominator below $10^{-6}$) are excluded, as standard. The chemical
potential $\mu$ is placed in the largest spectral gap in the filling window
$[0.02,0.30]$; here $\mu=-11.126$, occupation $n_{\text{occ}}=40$. The topological
contribution is isolated by subtracting a \emph{uniform ferromagnet} reference
($\mathbf n_i=\hat z$) at matched filling. Units follow the shared-kernel convention:
charge $\to e^2/h$, orbital $\to e/(2\pi)$.

\section{Results}
\begin{table}[h]
\centering
\begin{tabular}{lS}
\toprule
Quantity & {Value $[e/2\pi]$}\\
\midrule
$\sigma^{L_z}_{xy}$ (hopfion)              & 4330.3 \\
$\sigma^{L_z}_{xy}$ (uniform FM, matched filling) & 4055.4 \\
$\sigma^{L_z}_{xy}$ (topological residual) & 275.0 \\
\bottomrule
\end{tabular}
\caption{Real-space $12\times12\times6$, $m/t=7$, $\lambda=8$, $\mu=-11.126$,
$n_{\text{occ}}=40$, PBC, no SOC. Wall time $5.3$~s (numpy 2.3.5).}
\end{table}

The response is manifestly finite \emph{without} spin--orbit coupling, and it is
nonzero for the textured system---the qualitative mechanism the paper predicts. A
coarse reciprocal-space cross-check (\texttt{replication\_run\_fast.json}, $n_k=4$,
$8\times8\times4$) reproduces the expected antisymmetric energy dependence of
$\sigma_{xy}(E)$ about $E=0$ (values at $\pm E$ mirror), consistent with the paper's
Fig.~3 shape.

\section{Comparison to the headline claim}
\begin{itemize}
\item \textbf{Confirmed (qualitative):} a hopfion produces a finite 3D orbital Hall
conductivity $\sigma^{L_z}_{xy}$ at the Fermi level for $m=7t$ with no SOC. The
mechanism---local emergent field $\to$ orbital Berry curvature $\to$ transverse orbital
current---is reproduced. The textured value clearly differs from the collinear
reference.
\item \textbf{Not confirmed (quantitative):} the paper reports conductivities in
SI-like units $(\Omega\!\cdot\!\text{cm})^{-1}$ from a full reciprocal-space Kubo
integration with modern-theory orbital-magnetization corrections. Our real-space
$e/2\pi$ residual of $275$ is a difference of two large numbers ($4330-4055$) and is
not a clean, converged topological separation (see failure analysis).
\end{itemize}

\section{Critique / limits}
The uniform ferromagnet has full lattice symmetry and, without SOC, should give
$\sigma^{L_z}_{xy}\approx0$. That the reference instead returns $4055$ shows the
itinerant operator $L_z=\tfrac12(Xv_y-Yv_x)$ carries a spurious trivial contribution
under periodic boundary conditions (the position operators are ill-defined on a torus).
Switching to open boundaries shrinks both numbers by $\sim5\times$ (hopfion $44.7$,
FM $53.8$) but the FM reference is still not zero. The rigorous fix is the
modern-theory-of-orbital-magnetization current operator, a substantially larger build.
The present result therefore supports the \emph{existence} of the effect and its
SOC-free origin, but not a converged quantitative Hopf-index-resolved conductivity.

\section{Verdict}
\textbf{PARTIAL.} Qualitative 3D orbital Hall mechanism and finite orbital Hall
response without SOC are confirmed; strict quantitative topological separation is a
known hard limit shared with the sibling paper gobel2024.

\end{document}
